% Book trim and typography.
\usepackage[
  paperwidth=6in,
  paperheight=9in,
  inner=0.78in,
  outer=0.68in,
  top=0.72in,
  bottom=0.78in,
  headsep=0.20in,
  footskip=0.38in
]{geometry}
\usepackage{libertinus}
\usepackage[scaled=0.86]{sourcecodepro}
\usepackage{microtype}
\usepackage{xcolor}
\usepackage{enumitem}
\usepackage{scrlayer-scrpage}
\usepackage{mdframed}
\usepackage{needspace}
\usepackage{titlesec}
\usepackage{etoolbox}

% --- House microtypography: the craft a fine book is always set with. --------
% Even colour and optical margins: protrusion hangs punctuation past the text
% edge, expansion varies glyph widths a hair to avoid rivers. (Both default on
% under LuaLaTeX; stated here so the intent is legible and final.)
\microtypesetup{protrusion=true,expansion=true,final=true}

% Ligatures follow the word, not the type case: suppress fi/ff/fl and kin across
% a morpheme seam (shelf|ful, not shelfful welded by an ff-ligature).
\usepackage{selnolig}

% No widows, no orphans; no paragraph ending on a hyphen; no two hyphenated
% lines stacked; no hyphen carried across a page turn.
\clubpenalty=10000
\widowpenalty=10000
\displaywidowpenalty=10000
\brokenpenalty=10000
\finalhyphendemerits=10000
\doublehyphendemerits=10000
% No lone word on a paragraph's last line (a "runt") -- the widow and orphan
% penalties above cannot catch it, because it is a break INSIDE the paragraph,
% not at a page edge, and the default \parfillskip (0pt plus 1fil) lets the last
% line stretch infinitely, so a one-word last line costs nothing. A finite
% \parfillskip makes a very short last line underfull instead, so the line
% breaker re-sets the paragraph to draw the word up. The threshold is low (the
% last line need only reach about a fifth of the measure), so an ordinary
% paragraph whose last line is already a few words is untouched.
\parfillskip=0pt plus 0.8\hsize
% At least three letters either side of a break, and hyphenate a shade less
% eagerly than the default.
\lefthyphenmin=3 \righthyphenmin=3
\hyphenpenalty=150

% Pandoc's hypersetup partial wraps pdfkeywords in \xmpquote, which only some
% hyperref/hyperxmp combinations define. Provide it if hyperref does not.
\AddToHook{package/hyperref/after}{%
  \providecommand{\xmpquote}[1]{#1}%
}

\definecolor{ink}{HTML}{171717}
\definecolor{muted}{HTML}{5C5C5C}
\definecolor{rulegray}{HTML}{C8C8C8}
\definecolor{papergray}{HTML}{F4F4F1}
\definecolor{accent}{HTML}{4B3E8A}
\definecolor{linkgreen}{HTML}{46621F}

\color{ink}
\setlength{\parindent}{1.1em}
\setlength{\parskip}{0pt}
\linespread{1.045}
% The label separation is stated, at the class's own value captured here, so
% a list nested inside a ledger entry (whose description sets labelsep=1.6em)
% does not inherit the ledger's wide hang; every other list is unchanged.
\newlength{\PressLabelSep}
\setlength{\PressLabelSep}{\labelsep}
\setlist[itemize]{leftmargin=1.25em,itemsep=0.18em,topsep=0.35em,labelsep=\PressLabelSep}
\setlist[enumerate]{leftmargin=1.35em,itemsep=0.18em,topsep=0.35em,labelsep=\PressLabelSep}

% A key-plus-prose table arrives from the table-shapes filter as a ledger: a
% description list inside PressLedger, set in the manpage form -- the key on a
% line of its own, the prose at full measure beneath it, hanging a little in.
% A two-column table squeezed the prose into half the measure; a ledger has no
% column to squeeze. The style is set inside the environment only, so an
% author's own definition lists elsewhere keep the LaTeX default. A zero-width
% label box (the 1.6em hang is all label separation) makes every key wider
% than its box, so nextline breaks after every key, even a two-letter one.
% pandoc's \tightlist would zero the space between entries, so it is a no-op
% here, and every entry asks for three lines before it starts, so a key is
% never left alone at the foot of a page. The entry's \item calls the plain
% \item the house copied ONCE in the preamble (\PressPlainItem, below, #286):
% a list nested in an entry runs the before hook again, and a copy taken there
% would copy the already-redefined \item into itself and loop.
\newenvironment{PressLedger}{%
  \let\tightlist\relax
  \setlist[description]{style=nextline,leftmargin=1.6em,labelindent=0pt,%
    labelwidth=0pt,labelsep=1.6em,%
    itemsep=0.45\baselineskip,parsep=0pt,topsep=0.5\baselineskip,%
    font=\normalfont,before=\PressLedgerNeedspace}}%
  {}
\newcommand*{\PressLedgerNeedspace}{%
  \renewcommand*{\item}{\Needspace{3\baselineskip}\PressPlainItem}}
% The ledger's caption sets as every house table caption does (the caption
% package setup below: small, italic, centred, no label), and the header row
% survives as a one-line legend in small capitals above the entries.
% \captionof runs \noindent and \sloppy for its paragraph and, being a
% command rather than an environment, leaves them set: bare, one ledger
% caption killed the house paragraph indent and loosened the spacing of every
% page after it (the #292 fault). Each macro is its own group, and the
% legend's \par sits inside the group so \small sets its baselines.
\newcommand{\PressLedgerCaption}[1]{{\captionof*{table}{#1}}}
\newcommand{\PressLedgerLegend}[1]{\par{\noindent\small\textsc{#1}\par}}
% A three-column ledger's middle cell (a line of filenames, say) sets ragged
% right on its own line under the key: justified, a short run of code spans
% stretched its spaces across the measure.
\newcommand{\PressLedgerSub}[1]{{\raggedright #1\par}}

\clearpairofpagestyles
\makeatletter
\ihead{\footnotesize\color{muted}\itshape\@title}
\makeatother
\ohead{\footnotesize\color{muted}\leftmark}
% The optional argument styles plain pages (chapter openers) as well:
% without it they carry no folio at all, and a short appendix becomes
% an entirely unnumbered section the table of contents points into.
\cfoot[\footnotesize\thepage]{\footnotesize\thepage}
\renewcommand*{\chapterpagestyle}{plain.scrheadings}
\automark[section]{chapter}
% An unnumbered chapter (a preface, an appendix written {.unnumbered}, a
% colophon) comes through as \chapter*, which sets no mark, so its pages
% carried the last numbered chapter's running head (#282). The
% figure-numbering filter follows each such heading with
% \PressUnnumberedChapter{<title>}, which sets the title as the mark with no
% number -- KOMA's own \addchapmark, the mark \addchap sets.
% An unnumbered chapter also contributes no chapter number to the List of
% Plates (#288): the figure counter runs on from the last numbered chapter,
% so a colophon's plate listed as "16.1". Its plates list by caption alone
% (caption's empty list format), and the filter's \PressNumberedChapter
% restores the numbered entries at the next numbered chapter.
\newcommand*{\PressUnnumberedChapter}[1]{%
  \addchapmark{#1}%
  \captionsetup{listformat=empty}%
}
\newcommand*{\PressNumberedChapter}{\captionsetup{listformat=default}}

\addtokomafont{chapter}{\normalfont\bfseries\color{ink}}
\addtokomafont{section}{\normalfont\bfseries\color{ink}}
\addtokomafont{subsection}{\normalfont\bfseries\color{ink}}
\addtokomafont{disposition}{\raggedright}
\RedeclareSectionCommand[beforeskip=0pt,afterskip=1.25\baselineskip]{chapter}
\RedeclareSectionCommand[beforeskip=1.6\baselineskip,afterskip=0.55\baselineskip]{section}
\RedeclareSectionCommand[beforeskip=1.2\baselineskip,afterskip=0.35\baselineskip]{subsection}

% The illustrations split into two lists. Unnumbered plates (the literary
% woodcut idiom) keep the Victorian "List of Plates" that KOMA's
% \listoffigures prints; the numbered informative figures
% (.figure/.chart/.map/.photo/.diagram) collect into a separate "List of
% Figures", written to a private .lof2 file by the figure-numbering Lua
% filter (\addcontentsline{lof2}) and printed here by \PressListOfFigures.
% The book's contents show chapters only (tocdepth 0), which would also
% suppress the depth-1 illustration lines, so each list raises the depth
% for itself (the toc-depth/List of Plates scar). Both list headings take
% the title case of "Contents" and "List of Figures" (#289).
\AddToHook{begindocument/end}{%
  \renewcommand*{\listfigurename}{List of Plates}%
  \addtocontents{lof}{\protect\setcounter{tocdepth}{1}}%
}
% \PressListOfFigures prints only when the manuscript declares numbered
% figures: the filter then sets \PressListOfFigurestrue (through the
% include-before slot, which the template sets ahead of the contents), and
% the list prints just before \mainmatter, so it sits in the roman front
% matter after the List of Plates. Emitted as the first body block it came
% after \mainmatter and took arabic folios 1 and 2 ahead of the Preface
% (#289). A book with no numbered figures never sets the flag and is
% unchanged. It reads the .lof2 the filter wrote, raising tocdepth for
% itself so the depth-1 lines survive the chapters-only contents.
\newif\ifPressListOfFigures
\makeatletter
\newcommand*{\PressListOfFigures}{%
  \begingroup
    \setcounter{tocdepth}{1}%
    \chapter*{List of Figures}%
    \@mkboth{}{}%
    \@starttoc{lof2}%
  \endgroup
}
\pretocmd{\mainmatter}{\ifPressListOfFigures\PressListOfFigures\fi}{}{%
  \PackageError{press}{could not patch \string\mainmatter\space for the List of
    Figures}{The house header prints the List of Figures just before
    \string\mainmatter; the patch failed, so the list would be silently
    dropped.}}
\makeatother

% Woodcut figures: no "Figure N" labels, quiet small italic captions,
% pinned where the source places them rather than floated.
\usepackage{graphicx}
\usepackage{float}
% The caption package is the single caption authority (loaded after float, the
% package it must cooperate with). It sets every caption small and italic with
% no automatic label -- the literary woodcut idiom shows no "Figure N:", and a
% numbered figure carries its own bold "Figure C.N." inside the caption text,
% set by the figure-numbering filter -- centred under the figure as before. It
% also DEFINES the unnumbered \caption* that the float package leaves broken, so
% a placed figure's caption no longer prints a stray asterisk above itself.
% hypcap=true carries the list/link anchor to the figure top, the job the
% now-removed hypcap package did (hypcap conflicts with caption and anchored at
% the caption, overshooting).
\usepackage{caption}
\captionsetup{font={small,it},labelformat=empty,labelsep=none,%
  justification=centering,singlelinecheck=true,hypcap=true}
% wrapfig runs text around a placed figure; the figure-numbering filter uses
% its parity-aware inner/outer sides (i/o) for place=wrap-inner/wrap-outer.
% The toolchain image ships it (texlive-latex-extra); where it is absent a
% wrap degrades to an ordinary figure so the book still builds.
\IfFileExists{wrapfig.sty}{\usepackage{wrapfig}}{%
  \newenvironment{wrapfigure}[2]{\begin{figure}[H]\centering}{\end{figure}}%
}
\floatplacement{figure}{H}
% A figure taller than the text block makes LuaLaTeX ship empty pages
% forever, silently. Rebind pandoc's image bounding with caption headroom:
% cap at \PressFigureCap tall, never wider than the line, never upscaled.
% Never raise the cap near \textheight. The cap defaults to the v1 6.3in
% here and is re-set per print profile (build/profile-geometry.tex), which
% is included after this header.
\providecommand{\PressFigureCap}{6.3in}
\newsavebox\PressImgBox
\makeatletter
% The height cap, applied through pandoc's \pandocbounded hook. pandoc defines
% that macro (and \pandoc@box) only when the manuscript still holds a native
% image at write time; a book whose every figure is replaced by the
% figure-numbering filter leaves it undefined, and the filter's own
% \pandocbounded calls would then be an undefined control sequence. So the
% else branch DEFINES \pandocbounded (over our own savebox) rather than only
% setting the graphicx keys -- the cap holds either way, and a book pandoc did
% define it for is byte-identical (the \ifdefined branch is unchanged).
\AtBeginDocument{%
  \ifdefined\pandocbounded
    \renewcommand*\pandocbounded[1]{%
      \sbox\pandoc@box{#1}%
      \Gscale@div\@tempa{\PressFigureCap}{\dimexpr\ht\pandoc@box+\dp\pandoc@box\relax}%
      \Gscale@div\@tempb{\linewidth}{\wd\pandoc@box}%
      \ifdim\@tempb\p@<\@tempa\p@\let\@tempa\@tempb\fi
      \ifdim\@tempa\p@<\p@\scalebox{\@tempa}{\usebox\pandoc@box}%
      \else\usebox{\pandoc@box}%
      \fi}%
  \else
    \providecommand{\maxwidth}{\linewidth}%
    \newcommand*\pandocbounded[1]{%
      \sbox\PressImgBox{#1}%
      \Gscale@div\@tempa{\PressFigureCap}{\dimexpr\ht\PressImgBox+\dp\PressImgBox\relax}%
      \Gscale@div\@tempb{\linewidth}{\wd\PressImgBox}%
      \ifdim\@tempb\p@<\@tempa\p@\let\@tempa\@tempb\fi
      \ifdim\@tempa\p@<\p@\scalebox{\@tempa}{\usebox\PressImgBox}%
      \else\usebox{\PressImgBox}%
      \fi}%
  \fi}
\makeatother

% Scene break: the staggered three-row asterisk array (Alice's "fairy dust"),
% emitted by the scene-break Lua filter when the aesthetic sets
% scene-break: fairy-dust. Four asterisks, three offset into the gaps, four --
% the outer rows share one grid (each star centred in a \PressFDStep box) and
% the middle row is kerned half a step so it lands between them. Defined but
% never called for any other book, so their output is byte-identical.
\providecommand{\textasteriskcentered}{\ensuremath{\ast}}
\newlength{\PressFDStep}\setlength{\PressFDStep}{1.9em}
\newcommand{\PressFDStar}{\makebox[\PressFDStep]{\textasteriskcentered}}
\newcommand{\PressFairyDust}{%
  \par\addvspace{1.5\baselineskip}%
  \centerline{\vbox{\offinterlineskip
    \hbox{\PressFDStar\PressFDStar\PressFDStar\PressFDStar}%
    \vskip0.35em
    \hbox{\kern0.5\PressFDStep\PressFDStar\PressFDStar\PressFDStar}%
    \vskip0.35em
    \hbox{\PressFDStar\PressFDStar\PressFDStar\PressFDStar}%
  }}%
  \par\addvspace{1.5\baselineskip}%
}

% Verse set-piece (::: verse via set-pieces.lua): the house stanza measure --
% inset from the text block, ragged right, upright (the poem is not emphasis),
% a stanza space above and below. Line breaks arrive as \\ from the filter.
\newenvironment{PressVerse}{%
  \par\addvspace{0.7em}\begingroup
  \leftskip=3em\rightskip=0pt plus 2em\parindent=0pt\noindent}%
  {\par\endgroup\addvspace{0.7em}}

% Render Markdown blockquotes as quiet, book-like callouts.
\renewenvironment{quote}
  {\begin{mdframed}[
      backgroundcolor=papergray,
      linecolor=rulegray,
      linewidth=0.4pt,
      leftmargin=0pt,
      rightmargin=0pt,
      innerleftmargin=0.8em,
      innerrightmargin=0.8em,
      innertopmargin=0.55em,
      innerbottommargin=0.55em,
      skipabove=0.8em,
      skipbelow=0.8em
    ]\small\itshape}
  {\end{mdframed}}

% Any other table is a grid. pandoc gives a short-source table `ll` columns,
% which never wrap and run past the text block; the table-shapes filter gives
% every column a relative width instead, so pandoc writes p{} columns that
% wrap, and brackets the table with these two macros: a size smaller and
% tighter gutters.
% Code in a grid cell keeps the house code-span rule (#287, below): never
% hyphenated, an existing hyphen a last-resort break. The table head's keep is
% the figure-numbering filter's sized \Needspace before every longtable (#286),
% so the grid brackets add only the size and the gutters.
\newcommand*{\PressGridBegin}{%
  \begingroup\small\setlength{\tabcolsep}{0.5em}}
\newcommand*{\PressGridEnd}{\endgroup}

% Keep headings from being stranded at page bottoms.
\pretocmd{\section}{\Needspace{5\baselineskip}}{}{}
\pretocmd{\subsection}{\Needspace{4\baselineskip}}{}{}
% Nor a definition term (#286): each \item of a description list asks for
% two lines, the term and the first line under it, so a term never sits at a
% page foot with its definition overleaf. Only an item of the description
% itself asks (\@currenvir is the innermost environment, so a list nested in
% a definition is untouched), and only a description list carries the rule.
% (A table head is kept with its first rows by the figure-numbering filter,
% which sizes a \Needspace from the table's own cells.) The plain \item is
% copied ONCE, here in the preamble: copied inside the hook, a description
% nested in a definition would copy the wrapper itself, and \item would call
% itself forever -- a silent hang at full CPU.
\NewCommandCopy{\PressPlainItem}{\item}
\makeatletter
\AddToHook{env/description/begin}{%
  \def\item{%
    \ifx\@currenvir\PressDescriptionEnv\Needspace{2\baselineskip}\fi
    \PressPlainItem}%
}
\makeatother
\newcommand*{\PressDescriptionEnv}{description}

% Constant-width text never hyphenates (#287). An identifier or a path broken
% at a line end gained a hyphen it does not have (SUP-/PORTED), and a span
% holding a real hyphen could break right after it (--paper-/warm), so the
% real character vanished into the break. Inside \texttt (pandoc's inline
% code) the span sets in the no-hyphenation language, and each real hyphen is
% boxed, so the line breaker finds no discretionary after it. (LuaTeX reads
% \exhyphenpenalty and \automatichyphenmode at the paragraph's end, so a
% setting scoped to the span cannot forbid that break; the box can.) So a
% long identifier is not left with no break at all and pushed into the
% margin, a span may break after a "/", "_", or "." with no glyph added --
% the character stays at the line end, so nothing is invented -- at a high
% penalty, so it is a last resort taken only when the line has no better
% break. A hyphen the span already has is a legal break too, the least bad
% place when nothing else fits: the boxed hyphen is followed by a break that
% adds no glyph (\PressCodeHyphenBreak), at a higher penalty than the "/",
% "_", "." points so TeX takes those first. Without it a run of hyphenated
% identifiers had no break at all and ran off the page. Spaces inside a span
% still break, prose hyphenation is unchanged, and code blocks are verbatim
% and never hyphenated.
\makeatletter
\NewCommandCopy{\PressPlainTT}{\texttt}
\newcommand*{\PressCodeSpan}{\ifdefined\l@nohyphenation\language\l@nohyphenation\fi}
\newcommand*{\PressCodeBreak}{\penalty3000\relax}
\newcommand*{\PressCodeHyphenBreak}{\penalty5000\relax}
\makeatother
\ExplSyntaxOn
\tl_new:N \l__press_code_tl
\DeclareRobustCommand*{\texttt}[1]{
  \tl_set:Nn \l__press_code_tl {#1}
  \tl_replace_all:Nnn \l__press_code_tl { - } { \mbox{-} \PressCodeHyphenBreak }
  \tl_replace_all:Nnn \l__press_code_tl { \_ } { \_ \PressCodeBreak }
  \tl_replace_all:Nnn \l__press_code_tl { / } { / \PressCodeBreak }
  \tl_replace_all:Nnn \l__press_code_tl { . } { . \PressCodeBreak }
  \PressPlainTT { \PressCodeSpan \l__press_code_tl }
}
\ExplSyntaxOff

% Better code blocks in narrow trim.
\usepackage{fvextra}
\DefineVerbatimEnvironment{Highlighting}{Verbatim}{breaklines,breakanywhere,commandchars=\\\{\}}
% pandoc uses Highlighting only for a fence with a recognised language tag; an
% untagged fence (or one tagged `text`) becomes plain verbatim, which does not
% wrap, so a line past the measure ran off the page edge (#279). Rebind
% verbatim to the same fvextra environment with the same break options, so
% both kinds of fence wrap identically. The remaining options make a fence
% that already fit render pixel-identical to the kernel verbatim it replaces
% (the v3 contract: a fix must not move a valid book's layout):
%   backgroundcolorboxoverlap=0pt  fvextra follows every line with
%       \vspace{-0.25pt} even with no background colour, which shortened the
%       pitch inside a block by 0.5pt a line; off, the strutted line boxes
%       stack at exactly \baselineskip, as the kernel's one-line paragraphs do.
%   tabsize=1  fancyvrb sets a tab as an 8-space box; the kernel sets it as
%       one space. pandoc expands tabs before LaTeX sees them, so this only
%       reaches a raw tex include, but it stays faithful there too.
%   listparameters  fancyvrb inserts \clubpenalty and \widowpenalty (both
%       10000 in this house) after a block's first line and before its last;
%       the kernel sets each line as its own paragraph with no such penalty,
%       so page breaks inside a block must stay where they were.
% The Highlighting environment above keeps fvextra's defaults: changing its
% spacing would move every tagged fence, which is a design change.
% backgroundcolorboxoverlap exists only in the fvextra releases that
% added the overlap \vspace; the pinned toolchain image's fvextra (1.6.1)
% has neither the key nor the shortened pitch, and an unknown key is a fatal
% xkeyval error, so the option is passed only where fvextra defines it.
\ifcsname KV@FV@backgroundcolorboxoverlap\endcsname
\RecustomVerbatimEnvironment{verbatim}{Verbatim}{breaklines,breakanywhere,backgroundcolorboxoverlap=0pt,tabsize=1,listparameters={\clubpenalty=0 \widowpenalty=0 }}
\else
\RecustomVerbatimEnvironment{verbatim}{Verbatim}{breaklines,breakanywhere,tabsize=1,listparameters={\clubpenalty=0 \widowpenalty=0 }}
\fi

% Plain title page driven by document metadata (books that want the
% treatise front matter include tex/title-page.tex after this file).
\makeatletter
\renewcommand{\maketitle}{%
  \begin{titlepage}
    \raggedright
    \vspace*{0.8in}
    {\fontsize{25}{29}\selectfont\bfseries \@title\par}
    \vspace{0.35in}
    {\Large\color{muted} \@subtitle\par}
    \vfill
    {\large \@author\par}
    \vspace{0.12in}
    {\small\color{muted} \@date\par}
    \vspace{0.45in}
    {\small\color{muted} \directlua{tex.sprint(-2, os.getenv('BOOK_PUBLISHER') or '')}\par}
  \end{titlepage}
}
\makeatother
